\emph{Predictive modeling in finance is hampered by class imbalance: fraud, loan defaults, and bankruptcies constitute less than 5\% of observations in most real-world datasets. Existing generative methods -- CTGAN, TVAE, and TabSyn -- faithfully replicate these skewed distributions with no mechanism to adjust the class ratio at generation time. Direct diffusion on raw tabular features, as in TabDDPM, fails on mixed-type data due to incompatibility with one-hot encoded categorical features (KS $\approx$ 0.77).}

\emph{This work presents RE-TabSyn (Rare-Event Enhanced Tabular Synthesis), a framework that applies Classifier-Free Guidance, i.e., CFG, with latent diffusion to enable explicit, inference-time control over synthetic class proportions. RE-TabSyn combines three components: a Variational Autoencoder that encodes heterogeneous tabular features into a smooth continuous latent space; a Transformer-based Latent Diffusion Model that generates samples within that space; and a CFG mechanism that lets practitioners specify the target minority ratio through a single guidance scale parameter $w$, without retraining. To the best of the authors' knowledge, this is the first application of CFG to tabular data synthesis, confirmed through a systematic review of 151 papers across 10 categories.}

\emph{RE-TabSyn is evaluated on six financial benchmark datasets -- credit risk, income prediction, marketing response, loan performance, and corporate bankruptcy -- with minority ratios from 4.8\% to 44.5\%, across three random seeds. Three findings emerge. First, RE-TabSyn achieves approximately 50\% minority representation within 1.7\% of target on average, boosting minority rates by up to $10\times$. Second, classifiers trained on RE-TabSyn's balanced synthetic data reach a mean minority F1-score of 0.472, surpassing the real imbalanced-data baseline of 0.458, despite an acceptable fidelity trade-off (KS 0.171 vs.\ TabSyn's 0.109). Third, no systematic memorisation is observed, with mean Distance to Closest Record exceeding 1.0 across all datasets. The core findings have been accepted for presentation at the I2IT International Conference.}
